% !TEX root = ../main.tex

\section{Task B: Vary the Cycle Time}

\subsection{a) Calculations:}

Gegeben sind drei Zykluszeiten und ein konstanter Duty Cycle von \(D = 66.6\%\).

\begin{table}[h]
	\centering
	\begin{tabular}{|l|c|c|c|c|}
		\hline
		             & \textbf{LOW} & \textbf{MID} & \textbf{HIGH} & \textbf{Generic}               \\ \hline
		Cycle Time   & 1000 ms      & 500 ms       & 200 ms        & $t_{Cycle}$                    \\ \hline
		Duty Cycle   & 66.6\%       & 66.6\%       & 66.6\%        & $D$                            \\ \hline
		LED ON Time  & 666 ms       & 333 ms       & 133 ms        & $t_{on} = D \cdot t_{Cycle}$   \\ \hline
		LED OFF Time & 334 ms       & 167 ms       & 67 ms         & $t_{off} = t_{Cycle} - t_{on}$ \\ \hline
	\end{tabular}
\end{table}

\subsection{b) Implementation:}

Die Zykluszeit wird zwischen drei Stufen (LOW, MID, HIGH) variiert. Für jede Stufe werden Ein- und Ausschaltzeit aus der aktuellen Zykluszeit und dem festen Duty Cycle berechnet.

Die Werte werden in einem Array gespeichert und in einer Schleife durchlaufen. Für jede Stufe blinkt die LED genau 15 Mal, bevor zur nächsten Zykluszeit gewechselt wird.

\begin{lstlisting}
  /* USER CODE BEGIN 2 */
  const float tCycleLow = 1000;
  const float tCycleMid = 500;
  const float tCycleHigh = 200;
  const float dutyCycle = 0.666; 
  
  uint32_t OnTime;
  uint32_t OffTime; 

  float cycles[3] = {tCycleLow, tCycleMid, tCycleHigh};
  /* USER CODE END 2 */

  /* Infinite loop */
  /* USER CODE BEGIN WHILE */
  while (1)
  {
    /* USER CODE END WHILE */
    for(int level = 0; level < 3; level++)
    {
      float currentCycle = cycles[level];

      OnTime = (uint32_t)(currentCycle * dutyCycle);
      OffTime = (uint32_t)(currentCycle - OnTime);

      for(int cnt = 0; cnt < 15; cnt++)
      {
        GPIOJ->ODR |= SEGDP_Pin;
        HAL_Delay(OnTime);
        GPIOJ->ODR &= ~SEGDP_Pin;
        HAL_Delay(OffTime);
      }
    }
    /* USER CODE BEGIN 3 */
  }
  /* USER CODE END 3 */
\end{lstlisting}

\subsection{c) Discussion:}

Die Berechnung ist generisch und wiederverwendbar. Der Cast von \texttt{float} auf \texttt{uint32\_t} ist notwendig für \texttt{HAL\_Delay()}, führt aber zu kleinen Rundungsfehlern.

\texttt{HAL\_Delay()} bleibt blockierend.
